\documentclass[11pt]{article}
\usepackage[margin=1in]{geometry}
\usepackage{amsmath,amssymb,amsthm,mathtools,booktabs}
\usepackage[hidelinks]{hyperref}

\newtheorem{lemma}{Lemma}
\newtheorem{proposition}[lemma]{Proposition}
\newtheorem{corollary}[lemma]{Corollary}
\theoremstyle{remark}
\newtheorem{remark}[lemma]{Remark}
\theoremstyle{definition}
\newtheorem{assumption}{Assumption}

\newcommand{\E}{\mathbb E}
\newcommand{\R}{\mathbb R}
\newcommand{\norm}[1]{\lVert #1\rVert}
\newcommand{\zh}{\widehat z}
\newcommand{\zpar}{\widehat z_{\parallel}}
\newcommand{\zperp}{\widehat z_{\perp}}
\newcommand{\Pperp}{P_{\perp}}
\newcommand{\Proj}{\Pi}

\title{How gradient noise becomes generator bias in multilevel Picard,\\
and what a certified gradient retraction does:\\ a one-step analysis}
\author{Draft for the constraint-aware MLP project}
\date{2026-10-07}

\begin{document}
\maketitle

\begin{abstract}
We analyse a single generator evaluation inside a multilevel Picard (MLP) recursion for
semilinear parabolic PDEs whose solution has ridge structure.
For the centred terminal-block gradient estimator we show that the estimation noise
orthogonal to the ridge direction has trace exactly $(d-1)c_h/M$, with $c_h$
independent of the dimension $d$ (Lemma~\ref{lem:perp}).
Passed through a quadratic generator, this noise produces a deterministic bias equal to
$-\tfrac{\lambda}{2}(d-1)c_h/M$ (Proposition~\ref{prop:quad}); through a norm
generator it produces a bias of order $\sqrt{d/M}$ (Proposition~\ref{prop:norm}).
Retracting the gradient onto a PDE-certified set before the generator bounds the
generator error by a $d$-independent constant, and when the set lies in the ridge
direction it makes the law of the corrected generator value exactly independent of $d$
(Proposition~\ref{prop:cert}).
For generators that are linear in the gradient along the ridge direction, the orthogonal
noise never reaches the generator (Proposition~\ref{prop:linear}), which explains why
the certified retraction does not improve equal-cost accuracy in that class.
The analysis concerns one generator call (equivalently one Picard level); how deeper
levels amplify the effect is documented empirically, not proved here.
\end{abstract}

\section{Setting}

Consider $\partial_t u+\mathcal L u+f(t,x,u,z)=0$, $u(T,\cdot)=g$, with
$z=\sigma\nabla_x u$ for a scalar $\sigma>0$ and
$\mathcal L u=b\cdot\nabla u+\tfrac{\sigma^2}{2}\Delta u$ for a constant drift
$b\in\R^d$. Fix a unit vector $w\in\R^d$, write $\Pperp=I-ww^\top$, and use the following
assumptions.

\begin{assumption}[Ridge structure]\label{ass:ridge}
(i) $g(x)=\varphi(w\cdot x)$ with $\varphi:\R\to\R$ Lipschitz with constant $L_\varphi$;
(ii) $b=\beta w$ for some $\beta\in\R$;
(iii) the exact solution is a ridge function, $u^\star(t,x)=\psi(t,w\cdot x)$, so that
$z^\star(t,x)=\sigma\,\partial_s\psi(t,w\cdot x)\,w$.
\end{assumption}

Assumption~\ref{ass:ridge}(iii) holds whenever $f$ depends on $z$ only through
$\norm{z}$ or through $w\cdot z$ and the PDE has a unique solution in the relevant class
(translation invariance in the directions orthogonal to $w$). It covers the ridge
log-sum-exp HJB (P1), the norm-driver HJB (P4), the counterexample of
Hutzenthaler and Nguyen~\cite{HN2025} ($w=e_1$), and the viscous-Burgers family
$u^\star=\mathrm{sigmoid}(at+\sum_i x_i)$ (P2/P3, $w=d^{-1/2}\mathbf 1$).

\paragraph{The estimator.}
At a child point $(\tau,x)$ with $h=T-\tau>0$, MLP estimates the gradient by Monte Carlo.
We analyse the centred terminal block
\begin{equation}\label{eq:zhat}
\zh=\frac1M\sum_{k=1}^M \Delta_k\,\frac{G_k}{\sqrt h},\qquad
\Delta_k=g\bigl(x+\beta h\,w+\sigma\sqrt h\,G_k\bigr)-g(x),\qquad G_k\stackrel{\rm iid}{\sim}\mathcal N(0,I_d).
\end{equation}
This is the complete child gradient at the first Picard level and the terminal
contribution at every deeper level. It is also the noise source used in the lower-bound
argument of~\cite{HN2025}.
Write $s=w\cdot x$, $\xi_k=w\cdot G_k$, $G_{\perp,k}=\Pperp G_k$, so that
$\Delta_k=\varphi(s+\beta h+\sigma\sqrt h\,\xi_k)-\varphi(s)$ depends on $G_k$ only
through $\xi_k$. Define
\[
\zpar=w\cdot\zh,\qquad \zperp=\Pperp\zh,\qquad
Q_M=\frac1{Mh}\sum_{k=1}^M\Delta_k^2,\qquad c_h(s)=\E Q_M=\frac{\E\Delta_1^2}{h}.
\]

\section{The orthogonal noise}

\begin{lemma}[Orthogonal noise is isotropic, mean zero, and scales with $d-1$]\label{lem:perp}
Under Assumption~\ref{ass:ridge}(i)--(ii):
\begin{enumerate}
\item[(i)] $\E\zperp=0$ and $\E\zh=\bigl(\E\zpar\bigr)w$;
\item[(ii)] the law of $\zpar$ and of $Q_M$ does not depend on $d$;
\item[(iii)] conditionally on $(\xi_1,\dots,\xi_M)$,
$\zperp\sim\mathcal N\bigl(0,\tfrac{Q_M}{M}\Pperp\bigr)$;
\item[(iv)] $\E\norm{\zperp}^2=(d-1)\,c_h(s)/M$, and
$c_h(s)\le 2L_\varphi^2(\sigma^2+\beta^2h)$;
\item[(v)] $\E\norm{\zperp}=\E\bigl[\sqrt{Q_M}\bigr]\,\E\chi_{d-1}/\sqrt M
\ \ge\ \E\bigl[\sqrt{Q_M}\bigr]\,\dfrac{d-1}{\sqrt{(d+1)M}}$,
where $\chi_{d-1}$ denotes a chi random variable with $d-1$ degrees of freedom.
\end{enumerate}
\end{lemma}

\begin{proof}
Because $G_k$ is standard Gaussian and $w\perp\operatorname{range}\Pperp$, the scalar
$\xi_k$ and the vector $G_{\perp,k}\sim\mathcal N(0,\Pperp)$ are independent, and the
pairs are independent across $k$. Since $\Delta_k$ is a function of $\xi_k$,
$\zpar=(M\sqrt h)^{-1}\sum_k\Delta_k\xi_k$ and $Q_M$ are functions of
$(\xi_1,\dots,\xi_M)$, whose law does not involve $d$; this gives (ii).
Conditionally on the $\xi$'s, $\zperp=(M\sqrt h)^{-1}\sum_k\Delta_k G_{\perp,k}$ is a
linear combination of independent $\mathcal N(0,\Pperp)$ vectors with fixed
coefficients, hence Gaussian with mean zero and covariance
$(M^2h)^{-1}\sum_k\Delta_k^2\,\Pperp=(Q_M/M)\Pperp$; this gives (iii) and the first part
of (i). The second part of (i) follows from $\zh=\zpar w+\zperp$.
For (iv), $\E[\norm{\zperp}^2\mid\xi]=\operatorname{tr}\bigl((Q_M/M)\Pperp\bigr)
=(d-1)Q_M/M$; take expectations. The bound on $c_h$ follows from
$|\Delta_k|\le L_\varphi(|\beta|h+\sigma\sqrt h|\xi_k|)$ and $(a+b)^2\le 2a^2+2b^2$.
For (v), conditionally on $\xi$, $\norm{\zperp}=\sqrt{Q_M/M}\;\chi_{d-1}$ with
$\chi_{d-1}$ independent of $\xi$. For a nonnegative random variable $X$, Hölder's
inequality with exponents $3/2$ and $3$ applied to $X^2=X^{2/3}X^{4/3}$ gives
$\E X^2\le(\E X)^{2/3}(\E X^4)^{1/3}$, that is
$\E X\ge(\E X^2)^{3/2}(\E X^4)^{-1/2}$. With $\E\chi_k^2=k$ and
$\E\chi_k^4=k(k+2)$ this yields $\E\chi_k\ge k/\sqrt{k+2}$; take $k=d-1$.
\end{proof}

Under Assumption~\ref{ass:ridge}(iii) the exact gradient has no orthogonal component,
$\Pperp z^\star=0$. Every unit of $\E\norm{\zperp}^2$ is therefore pure estimation noise,
and it grows linearly in $d$ at fixed $M$.

\section{What a curved generator does with this noise}

\begin{proposition}[Quadratic generator: exact Jensen bias]\label{prop:quad}
Let $f(z)=-\tfrac\lambda2\norm{z}^2$, $\lambda>0$. Under Assumption~\ref{ass:ridge}(i)--(ii),
\[
\E f(\zh)-f(\E\zh)=-\frac\lambda2\Bigl(\operatorname{Var}(\zpar)+\frac{(d-1)\,c_h(s)}{M}\Bigr).
\]
If in addition Assumption~\ref{ass:ridge}(iii) holds, then
$\E f(\zh)-f(z^\star)=A_M(s)-\tfrac\lambda2(d-1)c_h(s)/M$, where
$A_M(s)=f(\E\zh)-f(z^\star)-\tfrac\lambda2\operatorname{Var}(\zpar)$ does not depend on $d$.
\end{proposition}

\begin{proof}
Since $\zperp\perp w$, $\norm{\zh}^2=\zpar^2+\norm{\zperp}^2$, so
$\E f(\zh)=-\tfrac\lambda2\bigl((\E\zpar)^2+\operatorname{Var}(\zpar)+\E\norm{\zperp}^2\bigr)$.
By Lemma~\ref{lem:perp}(i), $f(\E\zh)=-\tfrac\lambda2(\E\zpar)^2$, and (iv) gives the
result. Under (iii), $z^\star=(w\cdot z^\star)w$ and $w\cdot z^\star$ depends only on
$(\tau,s)$; together with Lemma~\ref{lem:perp}(ii) this makes $A_M$ independent of $d$.
\end{proof}

\begin{proposition}[Norm generator: bias of order $\sqrt{d/M}$]\label{prop:norm}
Let $f(z)=-\lambda\norm{z}$, $\lambda>0$. Under Assumption~\ref{ass:ridge}(i)--(ii),
\[
\E\bigl[\sqrt{Q_M}\bigr]\frac{d-1}{\sqrt{(d+1)M}}\ \le\ \E\norm{\zh}\ \le\
\sqrt{(\E\zpar)^2+\operatorname{Var}(\zpar)+\frac{(d-1)c_h(s)}{M}}.
\]
Consequently, under Assumption~\ref{ass:ridge}(iii),
$f(z^\star)-\E f(\zh)\ \ge\ \lambda\Bigl(\E\bigl[\sqrt{Q_M}\bigr]\tfrac{d-1}{\sqrt{(d+1)M}}-\norm{z^\star}\Bigr)$,
where $\E\sqrt{Q_M}$ and $\norm{z^\star}$ do not depend on $d$.
\end{proposition}

\begin{proof}
The lower bound uses $\norm{\zh}\ge\norm{\zperp}$ and Lemma~\ref{lem:perp}(v); the upper
bound is Jensen's inequality $\E\norm{\zh}\le(\E\norm{\zh}^2)^{1/2}$ with the
decomposition used in Proposition~\ref{prop:quad}.
\end{proof}

\begin{remark}[The Hutzenthaler--Nguyen counterexample]
In~\cite{HN2025}, $f_d(z)=\norm{\Pperp z}$ with $w=e_1$, $\sigma=1$, $\beta=0$ and
$g(x)=|x_1|$. Then $f_d(\zh)=\norm{\zperp}$ exactly, so
$\E f_d(\zh)^2=(d-1)c_h/M$ and $\E f_d(\zh)\ge \E[\sqrt{Q_M}]\,(d-1)/\sqrt{(d+1)M}$,
while $f_d(z^\star)=0$. At the origin $c_h=\E[(\sqrt h|\xi|)^2]/h=1$, which recovers the
identity $\E f_d(V_N)^2=(d-1)/N$ used in Appendix~B of the manuscript. The counterexample
is thus the special case in which the generator sees \emph{only} the orthogonal noise.
\end{remark}

\section{What a certified retraction does}

Let $\mathcal C(\tau,x)\subseteq\R^d$ be closed and convex with
$z^\star(\tau,x)\in\mathcal C(\tau,x)$ (a valid certificate), and let $f$ be
$L_{\mathcal C}$-Lipschitz on $\mathcal C$.

\begin{proposition}[Certified retraction caps the generator error]\label{prop:cert}
\begin{enumerate}
\item[(a)] For every realisation,
$\bigl|f(\Proj_{\mathcal C}\zh)-f(z^\star)\bigr|\le
L_{\mathcal C}\min\bigl(\norm{\zh-z^\star},\operatorname{diam}\mathcal C\bigr)$.
\item[(b)] If $\mathcal C\subseteq\operatorname{span}(w)$, write
$\mathcal C=\{cw:c\in I\}$ for a closed interval $I$ and let $\pi_I$ be the projection
onto $I$. Then $\Proj_{\mathcal C}\zh=\pi_I(\zpar)\,w$, so
$\norm{\Proj_{\mathcal C}\zh-z^\star}=|\pi_I(\zpar)-w\cdot z^\star|\le|\zpar-w\cdot z^\star|$.
Under Assumption~\ref{ass:ridge}, if $I$ depends only on $(\tau,w\cdot x)$ and
$f(cw)=F(c)$ for a function $F$ that does not depend on $d$, the law of
$f(\Proj_{\mathcal C}\zh)$ does not depend on $d$.
\end{enumerate}
\end{proposition}

\begin{proof}
(a) $z^\star=\Proj_{\mathcal C}z^\star$, both projections lie in $\mathcal C$, and metric
projection onto a closed convex set is nonexpansive; also
$\norm{\Proj_{\mathcal C}\zh-z^\star}\le\operatorname{diam}\mathcal C$.
(b) For $y\in\R^d$ and $c\in\mathcal C\subseteq\operatorname{span}(w)$,
$\norm{y-c}^2=\norm{\Pperp y}^2+\norm{(w\cdot y)w-c}^2$, so the minimiser over
$\mathcal C$ depends on $y$ only through $w\cdot y$; for $y=\zh$ it is
$\pi_I(\zpar)w$. The law statement follows from Lemma~\ref{lem:perp}(ii).
\end{proof}

\begin{remark}[Which component of the certificate does what]
Part (b) shows that the subspace component removes the $d$-dependence: the orthogonal
noise is discarded exactly. Part (a) shows that the magnitude component bounds the
generator error by $L_{\mathcal C}\operatorname{diam}\mathcal C$. Both are needed in a
recursion with an unbounded generator. Projecting onto $\operatorname{span}(w)$ alone
removes the $d$-growth but leaves the heavy tails of $\zpar^2$ to be amplified at deeper
levels; this matches the round-1 ablation on P1, where the span-only correction has mean
skill $573$ (median $2.35$) while the box certificate has $0.17$. For P1 the segment,
its box hull and the ball all have diameters independent of $d$, so (a) gives a
$d$-independent bound for each of them.
\end{remark}

\begin{corollary}[One averaged nonlinear correction]\label{cor:onelevel}
Fix a parent $(t,x)$. Let child times $\tau_i\in[t,T)$ be i.i.d., child positions
$X_i=x+\beta(\tau_i-t)w+\sigma W_{\tau_i-t}$, and let $\zh_i$ be the estimator
\eqref{eq:zhat} at $(\tau_i,X_i)$ with fresh samples. For nonnegative weights $a(\tau)$
consider $\widehat C=K^{-1}\sum_{i=1}^K a(\tau_i)f(\zh_i)$ and
$\widehat C_{\mathcal C}=K^{-1}\sum_{i=1}^K a(\tau_i)f(\Proj_{\mathcal C}\zh_i)$.
Under Assumption~\ref{ass:ridge}, with $C^\star=\E[a(\tau)f(z^\star(\tau,X))]$:
\begin{enumerate}
\item[(i)] for $f=-\tfrac\lambda2\norm z^2$,
$\E\widehat C-C^\star=B_0-\tfrac\lambda2\,(d-1)\,\E[a(\tau)c_{T-\tau}(w\cdot X)]/M$
with $B_0$ independent of $d$;
\item[(ii)] in the setting of Proposition~\ref{prop:cert}(b),
$|\E\widehat C_{\mathcal C}-C^\star|\le L_{\mathcal C}\,
\E\bigl[a(\tau)\min(|\zpar-w\cdot z^\star|,|I|)\bigr]$, a quantity independent of $d$.
\end{enumerate}
\end{corollary}

\begin{proof}
The law of $w\cdot X_i=s+\beta(\tau_i-t)+\sigma\,w\cdot W_{\tau_i-t}$ does not depend on
$d$. Condition on $(\tau_i,X_i)$ and apply Proposition~\ref{prop:quad} for (i) and
Proposition~\ref{prop:cert} for (ii).
\end{proof}

\section{The boundary: generators linear in the gradient along the ridge}

\begin{proposition}[Orthogonal noise is invisible to ridge-linear generators]\label{prop:linear}
Let $f(u,z)=A(u)+B(u)\,(v\cdot z)$ with $v\in\operatorname{span}(w)$. Then
$f(\widehat u,\zh)=A(\widehat u)+B(\widehat u)\,(v\cdot w)\,\zpar$ for any value
estimate $\widehat u$; in particular $f(\widehat u,\zh)$ does not depend on $\zperp$, and
conditionally on $\widehat u$ the generator is affine in $\zh$, so the gradient noise
produces no Jensen gap. Any remaining bias enters through the dependence between
$\widehat u$ and $\zpar$.
\end{proposition}

\begin{proof}
$v\cdot\zperp=0$ because $v\in\operatorname{span}(w)$ and $\zperp\perp w$.
\end{proof}

The viscous-Burgers family has $f(u,z)=\sigma u\sum_iz_i=\sigma\sqrt d\,u\,(w\cdot z)$
with $w=d^{-1/2}\mathbf 1$, so Proposition~\ref{prop:linear} applies. This is consistent
with three empirical findings: the certified box does not win at equal cost on P2/P3,
its benefit is confined to deep and under-sampled configurations, and the oracle
decomposition attributes part of the P2 error to the correlation between $u$- and
$z$-noise rather than to the gradient channel alone.

\section{Scope}

\begin{itemize}
\item The results are exact statements about one generator evaluation (and one averaged
correction) with the terminal-block estimator. At deeper Picard levels the child gradient
also contains nonlinear-correction terms; we do not prove how these propagate. The
empirical oracle decomposition and the $d$-sweeps in the mechanism suites are the
evidence for the recursive picture.
\item The ridge assumption is what makes the statements exact. It is also the regime in
which nonlinearity matters in high dimension for our benchmarks: when the terminal
condition depends on many directions, concentration of measure makes the nonlinear
correction nearly independent of $x$ (Rosenbrock HJB, multi-direction log-sum-exp).
\item All statements about the retraction require a valid certificate. Retractions that
are not the identity at $z^\star$ (illegal tighter sets, shrinkage towards a centre)
change the generator at the solution and are not covered.
\end{itemize}

\section{Numerical check}

The script \texttt{experiments/theory\_onestep/verify\_onestep\_theory.py} samples
\eqref{eq:zhat} by brute force in $\R^d$ (it does not use the conditional law of
Lemma~\ref{lem:perp}) at $s=0.3$, $h=0.2$, for the P1 terminal
$\varphi(s)=-\log\tfrac13\sum_k e^{\lambda_k s}$, $\lambda=(0.5,1.5,3)$,
$\sigma=\sqrt2$, $f=-\tfrac12\norm z^2$, certificate
$I=[-3\sqrt2,-0.5\sqrt2]$, and for the P4 terminal $\varphi(s)=\beta^{-1}\log\cosh(\beta s)$,
$\beta=2$, $\sigma=\sqrt2$, $f=-\norm z/\sqrt2$, certificate $I=[-\sqrt2,\sqrt2]$.
It uses $20{,}000$ replicates for $d\le200$ and $6{,}000$ for $d=1000$.
Full output, including $M=4$, is in
\texttt{results/theory\_onestep/onestep\_theory\_check.md}.
The ``certified'' columns use the projected generator minus $f$ evaluated at the
projected mean gradient of the linear problem, a common reference across $d$.

\begin{table}[h]\centering\small
\caption{Monte Carlo check at $M=16$. Quadratic rows: Jensen gap $\E f(\widehat z)-f(\E\widehat z)$; norm rows: $\E\lVert\widehat z\rVert$ with the bounds of Proposition~\ref{prop:norm}. Generator bias is measured against the common reference described in the text.}
\label{tab:check}
\begin{tabular}{@{}llrrrrrrr@{}}\toprule
driver & $d$ & \multicolumn{2}{c}{$\E\lVert\widehat z_\perp\rVert^2$} & \multicolumn{2}{c}{Jensen gap or $\E\lVert\widehat z\rVert$} & raw bias & \multicolumn{2}{c}{certified}\\
\cmidrule(lr){3-4}\cmidrule(lr){5-6}\cmidrule(lr){8-9}
 & & MC & theory & MC & theory / bounds & & bias & median\\\midrule
quadratic & 10 & 4.53 & 4.54 & -2.81 & -2.81 & -2.81 & -0.225 & 0.404\\
quadratic & 50 & 24.84 & 24.70 & -12.95 & -12.88 & -12.98 & -0.245 & 0.347\\
quadratic & 200 & 100.61 & 100.33 & -50.84 & -50.70 & -50.85 & -0.231 & 0.373\\
quadratic & 1000 & 501.66 & 503.66 & -251.34 & -252.34 & -251.31 & -0.199 & 0.397\\
norm & 10 & 0.49 & 0.48 & 0.845 & [0.607, 0.922] & -0.29 & -0.037 & 0.011\\
norm & 50 & 2.62 & 2.61 & 1.642 & [1.529, 1.725] & -0.86 & -0.032 & 0.014\\
norm & 200 & 10.69 & 10.62 & 3.195 & [3.137, 3.314] & -1.95 & -0.035 & 0.013\\
norm & 1000 & 53.31 & 53.31 & 7.058 & [7.029, 7.326] & -4.69 & -0.032 & 0.012\\
\bottomrule\end{tabular}\end{table}


In every row the orthogonal trace and the quadratic Jensen gap agree with
Lemma~\ref{lem:perp}(iv) and Proposition~\ref{prop:quad} within Monte Carlo error, the
norm expectation lies between the bounds of Proposition~\ref{prop:norm}, the raw
generator bias grows linearly (quadratic) or like $\sqrt d$ (norm), and the
distribution of the certified generator value is unchanged from $d=10$ to $d=1000$, as
Proposition~\ref{prop:cert}(b) predicts.

\begin{thebibliography}{1}
\bibitem{HN2025}
M.~Hutzenthaler and T.~A.~Nguyen.
\newblock Full history recursive multilevel Picard approximations suffer from the curse
of dimensionality for the Hamilton--Jacobi--Bellman equation of a stochastic control
problem.
\newblock arXiv:2506.23969, 2025.
\end{thebibliography}

\end{document}
